\section{相关工作}
\subsection{已观测旅行时间的分解}
稀疏探测车辆的两次位置观测可能跨越多个道路链路，因此需要沿已识别路径分配两次观测之间的耗时。Hellinga等利用链路几何与自由流速度，将旅行时间分为自由流、拥堵和停止时间\cite{hellinga2008}。其拥堵似然利用前一采样区间的信息，停止似然则反映车辆位置与链路下游端的关系，最终通过对可能的拥堵程度积分得到时间分配。Hofleitner和Bayen从干道排队与驾驶行为的统计模型推导链路时间分布\cite{hofleitner2011}。这些分布由对数凹分量混合而成，整体分配问题并非凸问题；固定延迟模式后可得到凸子问题，进而采用枚举或基于EM的算法求解。两项研究均要求分配时间之和等于观测总时长。

网络层面的推断还可以结合多次行程的共同观测。Yin等从行程进出记录推断链路时间分布，研究高斯似然、路径未知时的混合模型估计，以及面向更一般分布的行程拆分\cite{yin2015entryexit}。PR-LTTE交替利用行程距离与耗时恢复路径，并通过非负最小二乘估计链路时间\cite{sun2022prltte}。这些方法估计共享的道路链路量，部分方法还需恢复缺失路径。本文已知有序的有向公交站间区间，每个区间可以跨越多个道路链路；共享估计随可观测上下文变化，重新分配则随具体行程变化。NNLS-Ridge参照用于检验这一已知序列条件下的固定区间身份系数。

\subsection{旅行时间表示学习}
DeepTTE结合地理卷积与循环网络，同时学习整条路径和局部路径的旅行时间\cite{wang2018deeptte}。其局部标签来自GPS窗口的时间戳差，窗口之间可以重叠；整条路径时间由独立的注意力预测头输出。DutyTTE针对未提供路线时的旅行时间不确定性，将预测路径对齐与上下文相关的区间不确定性建模结合起来\cite{mao2025dutytte}。这类设计利用局部观测与路线表示预测尚未知晓的行程耗时。

ProbETA与重复链路的信息共享更为接近\cite{xu2025probeta}。Xu等将链路时间写为共享均值、日期随机效应和行程随机效应之和，通过学习低秩协方差表示描述行程之间和行程内部的依赖。训练利用行程总时间与时间戳生成的子行程时间构建联合似然，推断则结合邻近已完成行程，对查询行程的时间分布作条件估计。这是一种以概率模型组织共享链路证据的方式。本文已知查询行程的总时长，由有限单区间标签监督内部时间分配；行程修正的幅度相对于共享基础时间受到限制，并在本次行程内加和为零，其作用是调整分配比例。

\subsection{聚合监督与加和一致性}
聚合观测学习研究如何根据组级监督估计个体标签。Zhang等推导聚合观测的似然，并在条件独立假设下讨论由均值或总和监督进行回归\cite{zhang2020aggregate}。其中，聚合误差用于训练个体预测器。本文通过比例缩放使每次输出严格符合给定总量，因此最终输出的总和误差不提供学习信号。训练和模型选择分别使用训练集与验证集中的可见单区间标签。

预测协调从另一个角度处理加和一致性。MinT利用预测误差协方差结合不同聚合层级的预测，使协调后的预测满足加和关系\cite{wickramasuriya2019mint}。本文的总量是已观测输入，需要学习的是未知的区间份额，随后将其缩放到该总量。相应地，实验在相同输入、标签预算和最终缩放下比较“共享基础估计加行程修正”与直接区间预测，检验施加一致性约束后如何学习行程内部的时间分配。
